\documentclass[UTF8,fontset=none,11pt,a4paper]{ctexart}
\usepackage{ipara-handbook}
\usepackage{amsmath,booktabs,tabularx,hyperref}
\usetikzlibrary{arrows.meta,positioning,fit,calc}

\handbooksetup{OS-B04 VM × File × I/O}{408 · OS Internal Bridge}{File Page · Residency · Fault · Writeback}

\tikzset{
  vm_node/.style={draw=iparaDeep,fill=iparaSoft,rounded corners=1.5mm,align=center,minimum width=2.5cm,minimum height=0.88cm,font=\small},
  cache_node/.style={draw=iparaGreen,fill=white,rounded corners=1.5mm,align=center,minimum width=2.6cm,minimum height=0.88cm,font=\small},
  disk_node/.style={draw=iparaDeep,double,double distance=1.2pt,fill=white,rounded corners=1.5mm,align=center,minimum width=2.6cm,minimum height=0.88cm,font=\small},
  warn_node/.style={draw=iparaAccent,fill=white,rounded corners=1.5mm,align=center,minimum width=2.5cm,minimum height=0.88cm,font=\small},
  flow_edge/.style={-{Latex[length=2mm]},thick,draw=iparaDeep},
  wb_edge/.style={-{Latex[length=2mm]},thick,dashed,draw=iparaGreen}
}

\begin{document}
\handbooktitle{VM × File × I/O}{文件内容怎样成为可映射物理页面，并在缺页、回写和 I/O 完成中保持身份}{408 · OS-B04}

\begin{corebox}{Mother Interface}
\[
\boxed{
\begin{aligned}
\text{File Object + File Range} &\leftrightarrow \text{Cached/Resident Page State} \\
&\leftrightarrow \text{Process VA Mapping} \\
&\xleftrightarrow{\text{I/O request / completion}} \text{Persistent Storage State}
\end{aligned}
}
\]
本 Bridge 专门连接“文件系统（内容身份）、虚拟内存（地址空间映射与驻留）与 I/O 子系统（设备传输与回写）”的三方交接契约；页表硬件翻译归入 X-B02，外设中断与 DMA 硬件归入 X-B03。
\end{corebox}

\begin{methodbox}{先定义接口两侧的词}
\begin{itemize}
  \item \kw{File Object + File Range（文件对象 + 范围）}：文件系统对“这一页内容来自哪个持久对象的哪一段”的逻辑身份。经典 UNIX 可写成 inode/object + page index/offset，但真正 Owner 是 OS-06/07 的文件对象与 byte-to-block 关系。
  \item \kw{Page Cache / Cached File Page（文件页缓存）}：OS 为 file-backed 内容维护的内存副本/驻留状态。Linux 的 \code{address\_space}/folio/page 是工程表示；普通 buffered I/O 与 \code{mmap} 往往在这里汇合，但 direct I/O 等路径可以绕过部分缓存语义。
  \item \kw{File-backed Mapping（文件后备映射）}：进程地址空间中的某段 VA 与文件对象范围建立关系；建立 mapping 不代表内容已经驻留。
  \item \kw{Page-In / Fill（装入/填充）}：当所需文件内容当前没有可用内存副本时，VM/FS 需要产生存储请求并在完成后建立可用 resident state；是否由 DMA、interrupt 或 polling 完成属于 OS-05/X-B03 的下游机制。
  \item \kw{Dirty + Writeback（脏状态与写回）}：内存中的 file-backed 内容比持久副本更新时，系统必须记录“有未持久化修改”，并在策略或显式 durability 请求要求时产生写回；谁触发、何时批量、何时才算 durable 由 OS-05/06/07 分责。
\end{itemize}
\end{methodbox}

\begin{methodbox}{Owners 与职责边界}
\begin{itemize}
  \item \kw{左侧 Owner (OS-06/07 文件系统)：}拥有 Inode、目录路径解析、逻辑块号到物理块号（LBN $\to$ PBN）的映射。
  \item \kw{中间 Owner (OS-04 虚拟内存)：}拥有 VMA 范围、PTE 权限/有效位、缺页异常处理与物理页置换策略。
  \item \kw{右侧 Owner (OS-05 I/O 系统)：}拥有 I/O 请求队列、电梯调度算法（SCAN/C-SCAN）、设备驱动与 DMA 控制。
  \item \kw{本 Bridge Owns：}文件内容身份怎样与 cached/resident page state、进程 VA mapping 和 I/O request/completion 发生稳定 handoff；它不拥有任一侧内部的数据结构或回收/调度策略。
\end{itemize}
\end{methodbox}

\section{核心机制：同一内容要追踪三种“关系”，不是宣称三个身份完全等价}

一段 file-backed 内容跨层时至少要同时保留三类关系：
\begin{enumerate}
  \item \textbf{持久内容关系}：某个 File Object 的某个 byte/page range，决定数据在文件语义中“是谁”；
  \item \textbf{缓存/驻留关系}：该内容当前是否已有可用的内存副本、它是否 dirty/writeback，以及谁引用这份 resident state；
  \item \textbf{虚拟映射关系}：某个进程的 VA region/PTE 是否把这份内容暴露为可访问地址。
\end{enumerate}
同一个 file range 可以暂时没有 resident page，也可以被多个进程映射到不同 VA；一个 resident page 还可能正处 I/O 中。Bridge 的价值就是防止把\kw{内容身份、驻留状态、虚拟地址}压成一个对象。

\section{母模型：统一页缓存、缺页装入与回写流转拓扑}

\begin{center}
\begin{tikzpicture}[x=3.40cm,y=1.55cm]
  % Top: Process Virtual Addresses
  \node[vm_node] (va1) at (0, 0.7) {Process A\\[-1pt]\scriptsize $VA = \text{\code{0x7f4000}}$};
  \node[vm_node] (va2) at (0, -0.7) {Process B\\[-1pt]\scriptsize $VA = \text{\code{0x551000}}$};

  % Middle: Unified Page Cache
  \node[cache_node] (page) at (1.4, 0) {Unified Page Cache\\[-1pt]\scriptsize \code{Frame \#4096 (RAM)}};

  % Bottom: Block I/O Layer & Disk
  \node[vm_node] (bio) at (2.6, 0.7) {Block I/O Layer\\[-1pt]\scriptsize \code{submit\_bio()}};
  \node[disk_node] (disk) at (2.6, -0.7) {Disk Data Blocks\\[-1pt]\scriptsize LBN 1024 $\to$ PBN 8848};

  % Connections
  \draw[flow_edge] (va1) -- node[above,font=\scriptsize]{PTE} (page);
  \draw[flow_edge] (va2) -- node[below,font=\scriptsize]{PTE} (page);
  \draw[flow_edge] (page) -- node[above,font=\scriptsize]{Writeback} (bio);
  \draw[flow_edge] (bio) -- (disk);
  \draw[wb_edge] (disk) -- node[below,font=\scriptsize]{Page-In (DMA)} (page);
\end{tikzpicture}
\end{center}

\section{双向路径推导：缺页/缓存未命中与脏页写回}

\subsection{路径 A：File-backed Fault / Fill}
\begin{enumerate}
  \item 进程访问 file-backed VA，硬件发现当前翻译/权限状态不足以完成访问并进入 fault 慢路径；
  \item VM 根据 region/mapping 找到对应 File Object + File Range；
  \item 检查是否已有可复用的 cached/resident 内容：
  \begin{itemize}
    \item \textbf{Resident/Cache Hit}：建立或修复当前进程的映射关系，满足权限后重试原访问；
    \item \textbf{Content Not Resident}：文件系统把 file range 翻译成所需存储请求，OS-05 接管 request/driver/device progress；当前 task 若采用阻塞路径，则通过 OS-B01 建立等待关系并让出 CPU；
  \end{itemize}
  \item I/O completion 先更新“目标内存内容已可用/请求已完成”等软件状态，再使相关等待者 Ready/Runnable；完成发现可以来自 interrupt、polling 或混合机制；
  \item task 以后真正被调度运行时，VM 完成必要 mapping/fault return，并重试原访问。
\end{enumerate}

\subsection{路径 B：Dirty File Page $\to$ Writeback}
\begin{enumerate}
  \item 某次写入使内存中的 file-backed 内容比其持久副本更新，系统把相应缓存状态记为 dirty；硬件 dirty bit 可以帮助发现写入，但不是所有 OS 的唯一 dirty Owner；
  \item 当 writeback policy、内存回收压力或显式 \code{msync/fsync} 一类接口要求推进持久状态时，FS/VM 形成对应文件范围的写回工作；
  \item OS-05 把它转换为设备 request 并完成传输；请求可能被合并、拆分、排队，也不要求“一页恰好一次 DMA”；
  \item completion 到达后，内核根据本次提交期间是否又发生新写入等状态，决定哪些内容可以视为不再 dirty/不再 writeback。若调用方请求的是 durability，还必须继续满足 OS-06/07 定义的持久化协议，不能把“设备请求完成”直接等同于“文件事务 durable”。
\end{enumerate}

\section{\texttt{read()/write()} vs \texttt{mmap()} 全景对比}

\begin{center}
\small
\renewcommand{\arraystretch}{1.25}
\begin{tabularx}{\linewidth}{>{\bfseries}p{0.18\linewidth} X X X}
\toprule
对比维度 & 经典 \code{read()} / \code{write()} & 内存映射 \code{mmap()} & 性能与机制本质 \\
\midrule
数据路径 & 通过系统调用显式把 file range 与用户 buffer 之间的数据推进 & 建立 VA 到 file-backed 内容的映射，之后访问由普通 load/store + fault 慢路径驱动 & 两者通常都可复用文件页缓存；是否额外复制、是否命中 cache 取决于具体路径与实现 \\
控制权入口 & 每次 read/write 调用都有 syscall 进入/返回路径；真正是否 block 还看数据是否立即可用 & 建立 mapping 需要 syscall，之后每次访问不必 syscall；首次/再次缺页仍会进入内核 & “系统调用次数更少”不自动推出总体更快，page fault、TLB、局部性和访问模式都进入成本 \\
地址空间状态 & 用户 buffer 与文件内容是两个显式对象 & 额外建立 file-backed VA region 与页表关系 & \code{mmap} 把文件访问变成地址空间语义，因此要同时承担 mapping/lifetime 成本 \\
共享语义 & 多进程共享取决于它们是否共享打开实例、共享内存或应用协议 & \code{MAP\_SHARED} 一类映射允许多个进程针对同一 file-backed 内容建立共享可见关系 & “可见”不等于并发安全，也不等于修改已持久化到设备 \\
\bottomrule
\end{tabularx}
\end{center}

\begin{warnbox}{Anti-Bridge：三个必须阻断的混淆}
\begin{itemize}
  \item \(\text{Page Fault} \neq \text{文件读取失败}\)：在文件映射区，首次缺页是按需加载数据的正常驱动机制，绝非错误。
  \item \(\text{Hardware Cache (L1/L2/L3)} \neq \text{OS Page Cache (RAM)}\)：前者由 CPU 硬件自动管理，存放物理内存数据副本；后者由操作系统软件用 DRAM 管理磁盘文件页。
  \item \(\text{mmap()} \text{ 返回成功} \neq \text{文件已读入内存}\)：\code{mmap()} 建立的是 file-backed VA mapping；具体页面可以按需驻留，也可能因预读、已有 cache 或显式 populate 等原因提前存在，不能把 mapping creation 与 residency 画等号。
\end{itemize}
\end{warnbox}

\begin{corebox}{30 秒极速调用协议}
做题遇到“mmap 读写、Page Cache 共享、脏页回写与 I/O 路径”时，严格按三层追踪：
\[
\boxed{\text{VA Mapping 对应哪个 File Range？}} \xrightarrow{\text{内容是否 Resident/Cached？}} \xrightarrow{\text{需要时提交 I/O 并等待 Completion}} \boxed{\text{Dirty 与 Durable 状态分别推进}}
\]
\end{corebox}

\end{document}
